Pour la minuterie, il faut d'abord trouvé la vitesse du participant après sa
collision avec le ballon, dans deux cas. Une collision plastique lorsqu'il est
attrapé par le participant, et une collision semi-plastique lorsqu'il n'est pas
attrapé. Hors, les équations du coefficients de restitution et de l'énergie
total du systèmes restent les mêmes pour les deux cas. Les voici, en fonction du
participant (p) et ballon (b)

\begin{align}
	e&=\frac{V_{pf}-V_{bf}}{V_{bi}-{V_{pi}}}\\
	m_pV_{pi}+m_bV_{bi}&=m_pV_{pf}+m_bV_{bf}
\end{align}

\subsection{Cas ou le ballon est attrapé}
Dans ce cas-ci, la collision est plastique. Donc le coefficient de restitution
est 0.
\begin{equation}
0=\frac{V_{pf}-V_{bf}}{V_{bi}-{V_{pi}}}
\end{equation}
Vue que la collision est plastique, le participant et le ballon sont obligé
d'avoir la même vitesse final. Le participant a attrapé le ballon après tout.
Donc avec l'équation de l'énergie posé au début de cette section;
\begin{align}
	m_pV_{pi}+m_bV_{bi}&=m_pV_{f}+m_bV_{f}\\
	m_pV_{pi}+m_bV_{bi}&=(m_p+m_b)V_{f}\\
	\frac{m_pV_{pi}+m_bV_{bi}}{m_p+m_b}&=V_{f}
\end{align}
Toutes les variables de cette équation sont connue, sauf la vitesse initial du participant, qui est donnée par le calcul de la trajectoire. Hors, une fois tout mis dans l'équation on obtient la vitesse final du participant et du ballon:
\begin{equation}
V_{pf}\approx5.6\text{m/s}
\end{equation}

\subsection{Cas ou le ballon n'est pas attrapé}
Le coefficient de restitution choisi est 0.8. Simplement car plus il est grand
plus le participant pert de la vitesse. Il n'est pas possible de simplifier
les vitesses comme dans le cas précédent. Hors, la démarche est identique,
mais la formule obtenue est plus lourde:
\begin{align}
m_p V_{pi}+m_b V_{bi}&=m_p V_{pf}+m_b V_{bf} \\
m_p V_{pi}+m_b V_{bi}&=m_p V_{pf}+m_b \cdot(-e(V_{bi}-V_{pi})+V_{pf}) \\
m_p V_{pi}+m_b V_{bi}&=m_p V_{pf}-m_b e(V_{bi}-V_{pi})+m_bV_{pf} \\
m_p V_{pi}+m_b V_{bi}+m_be(V_{bi}-V_{pi})&=m_pV_{pf}+m_bV_{pf} \\
m_p V_{pi}+m_b(V_{bi}+e(V_{bi}-V_{pi}))&=V_{pf}(m_p+m_b) \\
\frac{m_pV_{pi}+m_b(V_{bi}+e(V_{bi}-V_{pi}))}{m_p+m_b}&=V_{pf}
\end{align}
Hors, avec la vitesse initial du participant restand la même que le cas
précédant, et le reste des variables étant fournis par le guide étudiant, la
vitesse final du participant s'il n'attrape pas le ballon est la suivante:
\begin{equation}
v_{pf}\approx5.07\text{m/s}
\end{equation}

\subsection{La minuterie}
La vitesse du participant dans les deux cas, étant calculé, il ne reste que la
minuterie. Hors, la longueur de la trappe étant connue (3m), le temps pour la
est donnée par:
\begin{equation}
\frac{L_T}{V_{pf}}=\text{ secondes}
\end{equation}
Ce qui donne ceci pour les deux cas:
\begin{align}
	T_{\text{attrapé}}&\approx0.54\text{ secondes}\\
	T_{\text{pas attrapé}}&\approx0.59\text{ secondes}
\end{align}
Hors, une marge de sécurité de 0.02 secondes est demandé, ce qui donne les
temps sécuritaire suivants:
\begin{align}
	T_{\text{attrapé}}&\approx0.55\text{ secondes}\\
	T_{\text{pas attrapé}}&\approx0.57\text{ secondes}
\end{align}
Hors, afin d'obtenir la réponse la plus sécuritaire, le juste millieu des
deux temps est choisi, par l'équation suivante:
\begin{equation}
T_m = \frac{T_{ba}+T_{bpa}}{2}\approx0.56\text{ secondes}
\end{equation}
